%=====================================================================
\chapter{Social, Legal and Ethical Issues}\label{ch:legal}
%=====================================================================
\locator{6}{Part VI --- Reach, Law and Practice}

\begin{outcomes}
By the end of this chapter you will be able to:
\begin{tight}
  \item \textbf{Name} the Pakistani instruments that bind an online shop, and
        \textbf{state} what each one governs.
  \item \textbf{Explain} what makes an online contract enforceable, and
        \textbf{identify} where consumer law actually comes from here.
  \item \textbf{Recognise} a dark pattern, including ones this course produced
        by accident.
  \item \textbf{Assess} what personal data a shop holds, and \textbf{justify}
        keeping or discarding each kind.
  \item \textbf{Argue} a professional position on a request you are asked to
        implement.
\end{tight}
\end{outcomes}

\begin{prereq}
\chref{basket}: what the shop records about visitors who never buy.
\chref{discounts}: the totals pipeline and what moving a line does.
\chref{security}: PECA 2016, and what the shop exposed. \chref{orders}: what a
refund does and does not do.
\end{prereq}

\begin{keyterms}
electronic contract $\bullet$ Electronic Transactions Ordinance 2002 $\bullet$
PECA 2016 $\bullet$ E-Commerce Policy 2019 $\bullet$ consumer protection
$\bullet$ provincial jurisdiction $\bullet$ personal data $\bullet$ data
minimisation $\bullet$ retention $\bullet$ dark pattern $\bullet$ accessibility
$\bullet$ WCAG $\bullet$ professional responsibility
\end{keyterms}

\section{What Makes an Online Order a Contract}

A shop takes money from strangers it never meets, on the strength of a click.
The question of why that click binds anybody is older than the software and is
answered by one instrument.

\begin{definitionbox}[Electronic Transactions Ordinance 2002]
Pakistan's \term{Electronic Transactions Ordinance 2002} gives legal
recognition to electronic documents, records and signatures: a contract is not
void merely because it was made electronically, and an electronic record
satisfies a legal requirement for writing.

\smallskip
Without it, every order in this course would be unenforceable. With it, the
click is a contract --- which cuts both ways, because it is also what makes
your shop's published terms binding on \emph{you}.

\smallskip
\textbf{Practical consequence:} the order record is the contract. That is why
\chref{checkout} copies the price, the address and the totals into the order
rather than pointing at the catalogue, and why \chref{orders} treats the order
history as the thing a dispute is settled from. Those were engineering
decisions in those chapters; here they are the reason they were the right ones.
\end{definitionbox}

\section{The Instruments, and What Each Governs}

\rowcolors{2}{white}{lightbg}
\begin{longtable}{@{}L{4.6cm}L{4.4cm}L{5.2cm}@{}}
\toprule
\rowcolor{primary!15}
\textbf{Instrument} & \textbf{Governs} & \textbf{Why it reaches you} \\
\midrule
\endhead
Electronic Transactions Ordinance 2002 & Electronic records, signatures,
contracts & Your orders are enforceable \\
\textbf{PECA 2016} & Offences involving information systems & Unauthorised
access is criminal --- \chref{security} \\
E-Commerce Policy of Pakistan 2019 & Policy direction: consumer protection,
payments, taxation, data & The framework the rest is built under \\
Provincial consumer protection Acts & Goods, services, redress & \textbf{Not
federal} --- see below \\
Sales tax law, administered by FBR & Registration, collection, returns &
\chref{shipping}'s rate and rules \\
\bottomrule
\end{longtable}

\begin{alertbox}[Two things about Pakistani law that surprise students]
\textbf{First: consumer protection is provincial, not federal.} Punjab, Sindh,
Khyber Pakhtunkhwa and Balochistan each have their own consumer protection
legislation, with their own authorities and their own procedures. A shop in
Bahawalpur selling to Karachi may be answerable under a different Act from the
one it was set up under, and ``the consumer protection law'' is not a thing
that exists.

\smallskip
\textbf{Second: Pakistan has no comprehensive data protection Act in force.} A
Personal Data Protection Bill has been in draft for years. What this means in
practice is \textbf{not} that there are no obligations --- contract, the
E-Commerce Policy, sectoral rules from the State Bank, and the plain law of
confidence all reach --- but that \emph{the law is not what will stop you doing
something stupid with a customer list}. Your own judgement is, and that is an
uncomfortable position to be in rather than a comfortable one.

\smallskip
\textbf{Both of these can change, and the second is likely to.} Check the
current position before advising anybody. What this chapter teaches is how to
ask, not what the answer was in 2026.
\end{alertbox}

\section{What Your Shop Already Knows About People}

Before any question of obligation, a question of fact: what is in there? This
course has already produced the answer, in pieces.

\begin{examplebox}[An inventory, from this course's own measurements]
\begin{tight}
  \item \textbf{645 session rows on a shop with no customers}
        (\chref{basket}) --- addresses, timestamps and a currency, for visitors
        who were mostly robots, kept for 24 hours by a clock rather than
        discarded on departure.
  \item \textbf{An address, a telephone number and an email} on every order,
        copied and kept permanently, for guests who never created an account
        (\chref{checkout}).
  \item \textbf{A complete purchase history per named person}, reachable in
        two clicks from the admin panel (\chref{orders}).
  \item \textbf{An IP address on every customer row}, written at registration
        and never updated or removed.
  \item \textbf{A PHP error log, readable by anyone over HTTP} before
        \chref{security}'s hardening --- and a PHP error log can contain the
        contents of the request that failed, which on a checkout is somebody's
        address.
  \item \textbf{Uploaded files in a web-served directory}
        (\chref{variations}), protected by an unguessable name.
\end{tight}

\smallskip
\textbf{Not one of those was a decision.} Every one is a default. The shop
collects them because the software does, and it keeps them because nothing
deletes them.
\end{examplebox}

\begin{conceptbox}[Three questions worth asking of each]
\begin{tightnum}
  \item \textbf{Why do we have this?} If the answer is ``it came with the
        software'', that is not an answer, and the thing to find out is what
        breaks if it goes.
  \item \textbf{How long do we need it?} An order is an accounting record and
        has a statutory life. A session row from a crawler has none. They are
        not the same object and should not have the same fate.
  \item \textbf{What happens if it leaks?} Rank by that, not by volume. The
        session table is large and dull; the order table is small and is
        somebody's home address next to what they bought.
\end{tightnum}

\smallskip
\textbf{Data minimisation is not a legal doctrine you are importing. It is the
engineering observation that you cannot lose what you did not keep}, and it is
free at design time and expensive afterwards.
\end{conceptbox}

\section{Dark Patterns, Including the Ones We Built}

\begin{definitionbox}[Dark pattern]
A \term{dark pattern} is an interface that leads a user to do something they did
not intend and would not have chosen, by exploiting how they read rather than
by lying to them. Each element is usually true. The arrangement is the
deception.
\end{definitionbox}

\begin{alertbox}[Three from this handout, two of them accidental]
\textbf{1. Charging tax on a discount the customer never received.}
\chref{discounts} measured it: moving the coupon line after the tax line
changed the customer's bill from 770.16 to 784.20. Every line was correct and
the total added up. \textbf{Nobody designed this as a deception --- it is one
integer in an admin field} --- and a shop that did it deliberately would be
indistinguishable from a shop that did it by accident. That is precisely why
the arithmetic in \chref{shipping} and \chref{discounts} is somebody's
professional responsibility and not a detail.

\smallskip
\textbf{2. Removing the confirmation screen to raise conversion.}
\chref{checkout} measured it as 28\% of the checkout's time, and listed it as
the one step not to remove. Removing it does raise completed orders. It raises
them by taking the last moment at which a customer sees what they are agreeing
to pay.

\smallskip
\textbf{3. Optimising a cash-on-delivery checkout for completion.}
\chref{payments} showed that an unverified COD order is a liability, and that
confirming orders by telephone is the most profitable step in the pipeline.
Advice imported from prepaid markets --- remove every step, reduce friction ---
produces more unwanted orders, and the customer who refuses at the door was
helped into an order they did not want.
\end{alertbox}

\begin{conceptbox}[A test you can apply before writing the code]
\textbf{Would you be comfortable explaining this arrangement to the customer it
affects, in one sentence, out loud?}

\begin{tight}
  \item ``We put the discount before the tax so you pay tax only on what you
        actually pay.'' Comfortable.
  \item ``We put the discount after the tax so you pay us an extra 1.8\%.''
        Not comfortable. Same software, one integer apart.
  \item ``We show you the total before you pay.'' Comfortable.
  \item ``We removed that screen because fewer people stopped.'' Not.
\end{tight}

\smallskip
The test is not a substitute for law and does not need to be. It is faster than
law, it applies before the code is written, and it has never once in this
course disagreed with the right answer.
\end{conceptbox}

\section{Accessibility, Which Is Not Optional and Is Usually Skipped}

\begin{conceptbox}[Four things that cost nothing at build time]
\chref{ux} rebuilt a theme without once mentioning the customers who cannot see
it. That omission is normal and is a mistake.

\begin{tightnum}
  \item \textbf{Every image has alt text}, and a product photograph's alt text
        is the product's name --- which the shop already has.
  \item \textbf{Every form field has a label} tied to it, not a placeholder
        pretending to be one. A placeholder disappears when typing starts,
        which is exactly when a confused customer needs it.
  \item \textbf{Colour is never the only signal.} ``Out of stock'' in red is
        invisible to a large minority of men.
  \item \textbf{The checkout works from the keyboard alone.} Try it. It takes
        four minutes and it is the only one of these four that usually fails.
\end{tightnum}

\smallskip
\textbf{All four are cheaper than page weight was in \chref{ux}}, and unlike
page weight they cannot be retrofitted by minifying anything. A shop that
cannot be used by somebody with a screen reader has excluded a customer who had
money and wanted to spend it, which is a commercial fact before it is an
ethical one.
\end{conceptbox}

\section{Being Asked to Do Something You Should Not}

\begin{examplebox}[Four requests you will receive, and an answer to each]
\begin{tight}
  \item \textbf{``Store the card number so customers need not retype it.''}
        No --- and the reason is not squeamishness: it moves the shop from PCI
        DSS SAQ-A to SAQ-D (\chref{gateway}), which is hundreds of
        requirements and an audit. \emph{Offer instead:} the provider's own
        tokenisation, which gives the same convenience and keeps the number
        where it is already somebody's job to protect.
  \item \textbf{``Email our customer list to this marketing company.''} Ask
        what the customers were told when they gave it, and whether that
        covers this. If nobody knows, the answer is no. \emph{Offer instead:}
        send the campaign yourself, which keeps the list where the promise
        was made.
  \item \textbf{``Delete these orders, they're embarrassing.''} No. They are
        accounting records and \chref{orders} showed they are also what a
        dispute is settled from. \emph{Offer instead:} a status, a note, and a
        report that excludes them --- which is what was actually wanted.
  \item \textbf{``Check whether our competitor's site has the same
        vulnerability.''} No, and say why plainly: that is unauthorised access
        to an information system, it is an offence under \textbf{PECA 2016},
        and no commercial reason changes that. \emph{Offer instead:} check
        your own, which is what \chref{security} did.
\end{tight}

\smallskip
\textbf{In every case the answer has the same shape:} say no once and clearly,
give the reason in terms of consequence rather than of principle, and offer the
nearest thing you can do. A refusal with no alternative gets overruled by
somebody who will not refuse.
\end{examplebox}

\begin{pitfall}
\begin{tight}
  \item \textbf{Assuming there is one national consumer protection law.} There
        is not; it is provincial.
  \item \textbf{Assuming no data protection Act means no obligations.} It means
        the law will not stop you, which is worse.
  \item \textbf{Treating a legal rate or rule as permanent.} Check it against
        the current instrument before advising anybody.
  \item \textbf{Keeping data because the software collected it.} That is not a
        reason, and it is not a defence.
  \item \textbf{Ranking data risk by volume} rather than by what a leak would
        do.
  \item \textbf{Building a dark pattern by accident}, which is what moving one
        integer did in \chref{discounts}.
  \item \textbf{Leaving accessibility until after launch}, when it cannot be
        retrofitted cheaply.
  \item \textbf{Refusing a request without offering an alternative.}
  \item \textbf{Testing any system you do not own} (\chref{security}).
\end{tight}
\end{pitfall}

\begin{lab}[Audit your own shop as though it were somebody else's]
\begin{lstlisting}[language=bash]
# --- 1. Inventory what you hold about people -------------------------
cd code
docker compose exec db mariadb -ushop -pshoppw opencart -e \
  "SELECT COUNT(*) FROM oc_customer; \
   SELECT COUNT(*) FROM oc_session; \
   SELECT COUNT(*) FROM oc_order; \
   SELECT COUNT(*) FROM oc_customer_ip;"

# --- 2. Find the most sensitive single row --------------------------
# Not the biggest table. The one whose leak would do most harm.
# Name it and say why in one sentence.

# --- 3. Decide a retention period for each --------------------------
# Three columns: what, how long, why. "Forever" needs a reason too.

# --- 4. Build the dark pattern, then look at it ---------------------
python m12_discounts.py --sortorder
# Read the two totals. Write the sentence you would have to say to
# the customer to explain the second one.

# --- 5. Put it back -------------------------------------------------
python m12_discounts.py --restore

# --- 6. Try your checkout with the keyboard only --------------------
# No mouse, no trackpad. Tab and Enter. Get to a completed order.
# Write down where you got stuck.

# --- 7. Check your images -------------------------------------------
curl -s "http://localhost:8090/index.php?route=product/category&path=20" \
  | grep -o '<img[^>]*>' | grep -c 'alt='
curl -s "http://localhost:8090/index.php?route=product/category&path=20" \
  | grep -c '<img'
# Two numbers. They should be equal.

# --- 8. Write your terms --------------------------------------------
# One page a customer could read: who you are, what you sell, the
# price including tax, delivery, how to cancel, how to complain, and
# what you do with their data. Then check your shop does what it says.
\end{lstlisting}

\smallskip
\textbf{What to notice.} Step 8 is the exercise. Almost every shop's published
terms describe a shop that does not exist, because nobody checked the sentences
against the software --- and this course has spent twenty-one chapters learning
how to check.
\end{lab}

\begin{checkpoint}
\begin{tightnum}
  \item A shop moves its coupon line after its tax line and collects 1.8\% more
        on every discounted order. Nothing was concealed and every figure is
        correct. Argue whether this is a dark pattern, state what would change
        your answer, and say what you would do if you found it in a shop you
        had inherited.
  \item Your shop holds 645 session rows and 6 orders. Rank the two tables by
        risk, justify the ranking, give a retention period for each with your
        reasoning, and say which one you would be more reluctant to delete and
        why.
  \item A client asks you to store card numbers ``because our competitor
        does''. Write the reply you would send: the refusal, the reason in
        terms of consequence, the alternative, and what you would do if the
        client insisted.
\end{tightnum}
\end{checkpoint}

\begin{chaptersummary}
\begin{tight}
  \item The Electronic Transactions Ordinance 2002 is what makes an online
        order a contract --- and makes your published terms binding on you.
        The order record is the contract, which is why it copies its data.
  \item PECA 2016 makes unauthorised access an offence; the E-Commerce Policy
        2019 sets the framework; FBR administers the tax.
  \item \textbf{Consumer protection in Pakistan is provincial}, with separate
        Acts and authorities, so ``the consumer protection law'' does not
        exist as a single thing.
  \item \textbf{Pakistan has no comprehensive data protection Act in force.}
        That means the law will not stop you, not that nothing binds you ---
        an uncomfortable position rather than a comfortable one.
  \item This course's own measurements are the data inventory: 645 session rows
        from robots, addresses and telephone numbers on guest orders, a
        purchase history per named person, an IP on every customer, and ---
        before hardening --- an error log anyone could read.
  \item Ask of each: why do we have it, how long do we need it, and what happens
        if it leaks. Rank by harm, not volume. You cannot lose what you did not
        keep.
  \item A dark pattern is an arrangement, not a lie, and \chref{discounts}
        built one by moving a single integer. Test by whether you could explain
        the arrangement out loud to the customer it affects.
  \item Four accessibility measures cost nothing at build time and cannot be
        retrofitted cheaply. The keyboard-only checkout is the one that usually
        fails.
  \item Refuse clearly, give the consequence rather than the principle, and
        always offer the nearest thing you can do.
\end{tight}
\end{chaptersummary}

\begin{reviewq}
  \item \textbf{(MCQ)} Consumer protection legislation in Pakistan is:
        (a)~federal only; (b)~provincial; (c)~governed by PECA 2016;
        (d)~not yet enacted anywhere.
  \item \textbf{(MCQ)} What makes an electronic order legally enforceable in
        Pakistan? (a)~PECA 2016; (b)~the E-Commerce Policy 2019; (c)~the
        Electronic Transactions Ordinance 2002; (d)~the shop's terms.
  \item \textbf{(Short)} Explain why an order copies its prices and addresses
        rather than referring to them, in legal rather than technical terms.
  \item \textbf{(Short)} Give the one-sentence test for a dark pattern, and
        apply it to a countdown timer that resets when the page is reloaded.
  \item \textbf{(Short)} Why is ``we have no data protection Act'' a reason for
        more care rather than less?
  \item \textbf{(Applied)} Write the data retention schedule for a shop taking
        50 orders a day: every table holding anything about a person, how long
        it is kept, the reason, what deletes it, and who checks. Mark anything
        you would need advice on rather than deciding yourself.
  \item \textbf{(Applied)} Your employer asks you to implement a checkout that
        pre-ticks an optional insurance product. Set out your analysis: what
        the law here does and does not say, what the one-sentence test gives,
        what you would propose instead, what you would do if overruled, and at
        what point you would decline to build it at all.
\end{reviewq}
